% This is the Reed College LaTeX thesis template. Most of the work
% for the document class was done by Sam Noble (SN), as well as this
% template. Later comments etc. by Ben Salzberg (BTS). Additional
% restructuring and APA support by Jess Youngberg (JY).
% Your comments and suggestions are more than welcome; please email
% them to cus@reed.edu
%
% See https://www.reed.edu/cis/help/LaTeX/index.html for help. There are a
% great bunch of help pages there, with notes on
% getting started, bibtex, etc. Go there and read it if you're not
% already familiar with LaTeX.
%
% Any line that starts with a percent symbol is a comment.
% They won't show up in the document, and are useful for notes
% to yourself and explaining commands.
% Commenting also removes a line from the document;
% very handy for troubleshooting problems. -BTS

% As far as I know, this follows the requirements laid out in
% the 2002-2003 Senior Handbook. Ask a librarian to check the
% document before binding. -SN

%%
%% Preamble
%%
\pdfobjcompresslevel 0 % pour être pris en charge par la plateforme d'archivage du CINES
% voir https://facile.cines.fr#latex
% \documentclass{<something>} must begin each LaTeX document
\documentclass[12pt,a4paper]{reedthesis}
% Packages are extensions to the basic LaTeX functions. Whatever you
% want to typeset, there is probably a package out there for it.
% Chemistry (chemtex), screenplays, you name it.
% Check out CTAN to see: https://www.ctan.org/
%%
\usepackage{graphicx,latexsym}
\usepackage{amsmath,amssymb,amsthm}
\usepackage{amsfonts}
\usepackage{longtable,booktabs,setspace}
\usepackage{chemarr} %% Useful for one reaction arrow, useless if you're not a chem major
\usepackage[hyphens]{url}
% Added by CII
\usepackage{lmodern}
\usepackage{float}
\floatplacement{figure}{H}
% End of CII addition
\usepackage{rotating}

\usepackage[utf8]{inputenc}
\usepackage[T1]{fontenc}
\usepackage{fancyhdr}
\usepackage{xcolor}
\definecolor{Prune}{RGB}{99,0,60}
\usepackage{mdframed}
\usepackage{multirow} %% Pour mettre un texte sur plusieurs rangées
\usepackage{multicol} %% Pour mettre un texte sur plusieurs colonnes
\usepackage{scrextend} %Forcer la 4eme  de couverture en page pair
\usepackage{tikz}
\usepackage[absolute]{textpos}
\usepackage{colortbl}
\usepackage{array}
%\RequirePackage{geometry}% That nicely create a one-page template
%\geometry{textheight=100ex,textwidth=40em,top=30pt,headheight=30pt,headsep=30pt,inner=80pt}
\usepackage{geometry}
\usepackage{hyperref}

% Next line commented out by CII
%%% \usepackage{natbib}
% Comment out the natbib line above and uncomment the following two lines to use the new
% biblatex-chicago style, for Chicago A. Also make some changes at the end where the
% bibliography is included.
%\usepackage{biblatex-chicago}
%\bibliography{thesis}


% Added by CII (Thanks, Hadley!)
% Use ref for internal links
\renewcommand{\hyperref}[2][???]{\autoref{#1}}
\def\chapterautorefname{Chapter}
\def\sectionautorefname{Section}
\def\subsectionautorefname{Subsection}
% End of CII addition

% Added by CII
\usepackage{caption}
\captionsetup{width=5in}
% End of CII addition

% \usepackage{times} % other fonts are available like times, bookman, charter, palatino

% Syntax highlighting #22


% Added by CII
%%% Copied from knitr
%% maxwidth is the original width if it's less than linewidth
%% otherwise use linewidth (to make sure the graphics do not exceed the margin)
\makeatletter
\def\maxwidth{ %
  \ifdim\Gin@nat@width>\linewidth
    \linewidth
  \else
    \Gin@nat@width
  \fi
}
\makeatother

%Added by @MyKo101, code provided by @GerbrichFerdinands

\renewcommand{\contentsname}{Table of Contents}
% End of CII addition

\setlength{\parskip}{0pt}

% Added by CII

\providecommand{\tightlist}{%
  \setlength{\itemsep}{0pt}\setlength{\parskip}{0pt}}

% End of CII addition
%%
%% End Preamble
%%
%

\begin{document}

\begin{titlepage}

%\thispagestyle{empty}

\newgeometry{left=7.5cm,bottom=2cm, top=1cm, right=1cm}
\tikz[remember picture,overlay] \node[opacity=1,inner sep=0pt] at (-28mm,-135mm){\includegraphics{logos/bandeau.pdf}};

% fonte sans empattement pour la page de titre
\fontfamily{fvs}\fontseries{m}\selectfont

%*****************************************************
%******** NUMÉRO D'ORDRE DE LA THÈSE À COMPLÉTER  ****
%******** après le premier dépôt légal /          ****
%******** French legal PhD number to be completed ****
%******** after the first legal deposit           ****
%*****************************************************

\color{white}
\begin{picture}(0,0)
\put(-150,-735){\rotatebox{90}{NNT~: 2020UPASA000}}
\end{picture}

%*************************************************************
%**  LOGO  ÉTABLISSEMENT PARTENAIRE SI COTUTELLE :          **
%**  CHANGER L'IMAGE logoCotutelle.png; SINON COMMENTER  /  **
%**  Logo of partner establishment if cotutelle agreement : **
%**  change image logoCotutelle.png; otherwise add % signs  **
%*************************************************************
\vspace{10mm}
\vspace{-20mm} % à ajuster en fonction de la hauteur du logo
\flushright \includegraphics[width=8.0cm]{logos/cotutelle.png}

%*****************************************************
%**************** TITRE / TITLE **********************
%*****************************************************
\flushright
\vspace{15mm} % largeur à régler éventuellement / width to adjust if necessary
\color{Prune}
\fontfamily{fvs}\fontseries{m}\fontsize{12}{16}\selectfont
  Étude de la germination et des réarrangements du protéome et du métabolisme pour assurer la qualité des graines et leur défense via la production de peptides: Approches intégratives et prédictives
\fontfamily{fvs}\fontseries{m}\fontsize{12}{16}\selectfont
  Proteome and metabolome remodeling during seed germinationInvestigation of germination and the rearrangement of the proteome and metabolism to enhance seed quality and defense through peptide production: Integrative and predictive approaches

%*****************************************************

%\fontfamily{fvs}\fontseries{m}\fontsize{8}{12}\selectfont
\normalsize
\vspace{15mm}

\color{black}
\textbf{Thèse de doctorat de l'Université Paris-Saclay}

\vspace{10mm}

\hspace*{-0.7cm}Ecole Doctorale numéro 581 : AGRICULTURE, ALIMENTATION, BIOLOGIE, ENVIRONNEMENT, SANTÉ (ABIES)\\
\small Spécialité de doctorat~: Bioinformatique\\
\footnotesize Unité de recherche~: Université Paris-Saclay, INRAE, AgroParisTech, Institut Jean-Pierre Bourgin, 78000, Versailles, France.\\
\footnotesize Référent~: AgroParisTech

\vspace{10mm}

\textbf{Thèse présentée et soutenue à Paris, le 01 2025, par}
\bigskip
\Large {\color{Prune} \textbf{Shuang PENG}}

\vspace{15mm}

\flushleft \normalsize \textbf{Composition du jury~:}
\bigskip

\scriptsize
\arrayrulecolor{Prune}
\begin{tabular}{|p{8cm}l}
\textbf{Béatrice TEULAT} &  Rapporteur 1\\
Professeure, L'Institut Agro Rennes-Angers, IRHS (UMR 1345) & \\
\textbf{Frédéric CAPEL} &  Rapporteur 2\\
Chargé de recherche, INRAE, Clermont-Ferrand, UNH (UMR 1019) & \\
\textbf{Juliette LEYMARIE} &  Examinateur 1\\
Professeure, Université Paris Est Créteil, iEES Paris & \\
\textbf{Thomas BLEIN} &  Examinateur 2\\
Chargé de recherche, CNRS (Université Paris-Saclay), IPS2 (UMR 1403) & \\
\textbf{Thomas SIMONSON} &  Examinateur 3\\
Directeur de Recherche, CNRS (IPP)/École Polytechnique, Laboratoire de Biochimie, UMR 7654 & \\
\textbf{} &  \\
 & \\
\end{tabular}

\begin{tabular}{p{8cm}l}
& \\
\textbf{} &  \\
 & \\
\textbf{} &  \\
 & \\
\textbf{} &  \\
 & \\
\end{tabular}

\end{titlepage}

\addamargin % this add the margin back

\frontmatter % this stuff will be roman-numbered
\pagestyle{empty} % this removes page numbers from the frontmatter








\mainmatter % here the regular arabic numbering starts
\pagestyle{fancyplain} % turns page numbering back on




% 4eme de couverture
\ifthispageodd{}{\newpage\thispagestyle{empty}\null}
% Use the next command with TeX Live version ≥ 2020
%\Ifthispageodd{}{\newpage\thispagestyle{empty}\null}
\newpage
\thispagestyle{empty}
\newgeometry{top=1.5cm, bottom=1.25cm, left=2cm, right=2cm}
\fontfamily{rm}\selectfont

\lhead{}
\rhead{}
\rfoot{}
\cfoot{}
\lfoot{}

\noindent
%*****************************************************
%***** LOGO DE L'EDMH *********
%*****************************************************
\includegraphics[height=2.45cm]{logos/logo\_usp\_ABIES.png}
\vspace{1cm}
%*****************************************************

\begin{mdframed}[linecolor=Prune,linewidth=1]
\vspace{-.25cm}
\paragraph*{Titre~:} Étude de la germination et des réarrangements du protéome et du métabolisme pour assurer la qualité des graines et leur défense via la production de peptides: Approches intégratives et prédictives
\begin{small}
\vspace{-.25cm}
\paragraph*{Mots clés~:} Sciences végétales, Semences, Germination, Protéomique, Peptides, Métabolomique.

\vspace{-.5cm}
\begin{multicols}{2}
\paragraph*{Résumé~:} La germination et l'émergence des plantules sont des étapes critiques pour les productions végétales. En effet, la phase de transition entre l'état hétérotrophe vers un développement autotrophe est extrêmement sensible aux contraintes environnementales. Ces travaux de thèse se concentrent sur l'analyse des descripteurs de germination, fournissant ainsi un outil précieux pour les physiologistes afin d'améliorer le phénotypage de la germination. Cette approche permettra également d'intégrer une grande diversité d'indices de germination dans des analyses génomiques. Les mécanismes de résistance/tolérance au stress thermique demeurent peu décrits à ces stades précoces de germination. Ce travail révèle les transitions métaboliques de la germination de la lentille sous l'effet de la température par des approches de protéomique et de métabolomique, et permet d'étudier les réorganisations du protéome et le potentiel de la germination pour améliorer l'intérêt nutritionnel des graines. Des approches d'annotation fonctionnelle de gènes et de modélisation statistique de données protéomiques et métabolomiques sur des graines en germination conduiront à mieux comprendre les mécanismes de réponse au stress et à caractériser des biomarqueurs de la qualité des semences. Enfin, ce travail de thèse consiste à mettre enouvre une approche d'intelligence artificielle pour prédire les peptides antimicrobiens générés par l'hydrolyse des protéines de réserve pendant la germination. L'ensemble des travaux développés dans cette thèse sera largement partagé avec la communauté scientifique du domaine des sciences végétales.
\end{multicols}
\end{small}
\end{mdframed}

\begin{mdframed}[linecolor=Prune,linewidth=1]
\vspace{-.25cm}
\paragraph*{Title:} Proteome and metabolome remodeling during seed germinationInvestigation of germination and the rearrangement of the proteome and metabolism to enhance seed quality and defense through peptide production: Integrative and predictive approaches

\begin{small}
\vspace{-.25cm}
\paragraph*{Keywords:} Plant science, Seed, Germination, Proteomic, Peptides, Metabolomic

\vspace{-.5cm}
\begin{multicols}{2}
\paragraph*{Abstract:} Germination and seedling emergence are critical stages for plant production. Indeed, the transition from a heterotrophic state to autotrophic development is highly sensitive to environmental constraints. This thesis work focuses on the analysis of germination descriptors, providing a valuable tool for physiologists to enhance germination phenotyping. This approach will also enable the integration of a wide range of germination indices into genomic analyses. Mechanisms of resistance/tolerance to thermal stress remain poorly described at these early stages of germination. This work unveils the metabolic transitions of lentil germination under the influence of temperature through proteomic and metabolomic approaches, allowing for the study of proteome rearrangements and the potential to enhance the nutritional value of seeds. Functional gene annotation approaches and statistical modeling of proteomic and metabolomic data on germinating seeds will lead to a better understanding of stress response mechanisms and the characterization of seed quality biomarkers. Finally, this thesis work involves implementing an artificial intelligence approach to predict antimicrobial peptides generated by reserve protein hydrolysis during germination. The entirety of the work developed in this thesis will be widely shared with the scientific community in plant sciences.
\end{multicols}
\end{small}
\end{mdframed}


\vfill
\fontfamily{fvs}\fontseries{m}\selectfont
\noindent\begin{tabular}{p{14cm}}
\multirow{3}{16cm}[+0mm]{\small {\color{Prune} {\bf Université Paris-Saclay}\\
{\scriptsize Espace Technologique / Immeuble Discovery}\\
{\scriptsize  Route de l’Orme aux Merisiers RD 128 / 91190 Saint-Aubin, France}}}\\\mbox{}
\end{tabular}


% Index?

\end{document}
